\centerline{\bf \Huge Регион Эйлера}

\begin{flushright}
        \scriptsize{}
\end{flushright}

\problem{1}
Можно ли белый квадратный лист бумаги разрезать на 25 квадратных частей и покрасить 4 из этих частей чёрным цветом так, чтобы стороны квадратов одного цвета были равны, а стороны квадратов разных цветов были различны? (Каждая из 25 частей должна быть либо чёрной, либо белой.)

\problem{2}
Дано натуральное число $n$, большее 3 . Все простые числа от 1 до $n$ включительно раскрасили в два цвета: синий и красный. Оказалось, что модуль разности произведения синих чисел и произведения красных чисел меньше $n$. Чему он может быть равен?

\problem{3}
Через вершину $D$ квадрата $A B C D$ проведена прямая $\ell$, не пересекающая сторон квадрата в других точках. Луч $C A$ пересекает $\ell$ в точке $E$. На продолжении отрезка $E D$ за точку $D$ выбрана точка $F$ так, что $\angle D C F=45^{\circ}$. Оказалось, что $A E=C F$. Найдите $\angle C F D$.

\problem{4}
Петя и Вася играют в следующую стратегическую игру. На столе лежат 100 карточек с натуральными числами от 1 до 100 (все числа разные) и 100 одинаковых фишек. Игра состоит из 100 раундов, каждый из которых заключается в следующем:

Сначала Петя выбирает одну из карточек со стола.

Затем Вася решает, что сделать с выбранной карточкой:

\quad
$-$ либо отдать её Пете и взять себе одну фишку со стола;

\quad
$-$ либо отдать одну из своих фишек и забрать карточку себе.

Изначально ни у Пети, ни у Васи нет ни карточек, ни фишек. Какую наибольшую сумму чисел на своих карточках Вася может гарантировать себе независимо от игры Пети?

\problem{5}
Каждая сторона равностороннего треугольника разделена на 6 равных частей, через точки деления проведены прямые, параллельные сторонам, делящие исходный треугольник на 36 маленьких треугольничков. В каждой из вершин этих треугольничков сидит по жуку. Они одновременно начинают двигаться по линиям деления с равными скоростями. Когда жук попадает в вершину треугольничка, он поворачивает на $60$ или $120$ градусов. Докажите, что через некоторое время какие-то два жука окажутся в одной вершине маленького треугольничка.